\documentclass[11pt]{article}
\usepackage[margin=1in]{geometry}
\usepackage{amsmath,amssymb,amsthm,hyperref}
\newcommand{\E}{\mathbb E}
\newcommand{\CN}{\mathrm{CN}}
\newcommand{\ph}{\operatorname{phase}}
\newtheorem{proposition}{Proposition}
\title{A symmetry-closed response map and original-noise replicas}
\author{Grothendieck certificate project}
\date{23 October 2026}
\begin{document}
\maketitle

\paragraph{Purpose and scope.}
The genuine cubic used to obstruct the frozen radial tables also supplies
actual bounded responses. This note retains the shared projection fields and
the original Gaussian noise coupling. It gives an exact finite-dimensional
image map for a restricted symmetry class and a fixed-size replica model for
its entire correlation kernel. It gives no global norm cap, improved numerical
endpoint, or complete upper-rounding score. The invariant-space and Gaussian
conditioning arguments are paper proofs. Seven replica algebra results and
four tail/price bounds are formalized separately in Lean.

\section*{1. The invariant first/cubic space}
Fix $m\ge2$, $d=m(m+1)/2$, and $D=m(m+1)$. Let $Z$, $X_i$, and $Y_{ij}$,
$1\le i\le j\le m$, be independent standard proper complex Gaussians. Put
\[
 q=\frac{\sum_i X_i^2\bar Y_{ii}
          +\sqrt2\sum_{i<j}X_iX_j\bar Y_{ij}}{\sqrt D},\qquad
 A_X=\frac{\|X\|^2-m}{\sqrt m},\qquad
 A_Y=\frac{\|Y\|^2-d}{\sqrt d}.
\]
Act by $U\in U(m)$ on $X$ and by its unitary symmetric-square representation
on $Y$, leaving $Z$ fixed. Then $q,A_X,A_Y$ are invariant. In addition require
charge one under simultaneous scalar rotation of all original Gaussian
coordinates. The first and $(2,1)$ cubic invariant modes are the orthonormal tuple
\begin{equation}\label{eq:five-modes}
 B=\left(Z,\frac{Z(|Z|^2-2)}{\sqrt2},ZA_X,ZA_Y,q\right).
\end{equation}

Here is an elementary classification. The central scalar subgroup gives weight
one to $X$ and weight two to $Y$. There is no invariant remote linear or
holomorphic quadratic. Among neutral quadratics, diagonal torus invariance,
coordinate permutations and two-coordinate rotations leave exactly the centered
norms of $X$ and $Y$. Among remote $(2,1)$ cubics, the central weights leave only
$XX\bar Y$. Diagonal torus invariance pairs each $X_iX_j$ with $\bar Y_{ij}$;
permutations and a two-coordinate rotation impose the relative coefficients
$1,\sqrt2$ displayed above. Combining remote and visible degrees yields
\eqref{eq:five-modes}. Wick ordering commutes with the unitary action.
The norms and orthogonality follow directly from independent Gaussian moments.

For $T=P_{1,0}-\epsilon P_{2,1}$, write
$\Lambda=\operatorname{diag}(1,-\epsilon,-\epsilon,-\epsilon,-\epsilon)$.
For $c\ne0$ define
\[
 f_c=\ph\left(\sum_{j=0}^4c_jB_j\right),\qquad
 \kappa_j(c)=\E[\bar B_j f_c].
\]
The exact best-response image is
\[
 Tf_c=\sum_jB_j(\Lambda\kappa(c))_j,
\]
and the pair payoff is
\begin{equation}\label{eq:finite-pair-payoff}
 J(c,e)=\Re\bigl(\kappa(c)^*\Lambda\kappa(e)\bigr).
\end{equation}
Thus the norm restricted to invariant, charge-one responses is the supremum
of \eqref{eq:finite-pair-payoff} over two five-component complex vectors.
Indeed the coefficient body of disc-valued responses is compact, since only
finitely many integrable probes are retained. A maximizing pair can be replaced
by mutual phase best responses without loss, and their images lie in
\eqref{eq:five-modes}. Conversely every displayed phase is an actual admissible
response. This is an exact restricted-class formulation; averaging an arbitrary
maximizer under $U(m)$ need not preserve its payoff, so no unrestricted equality
is asserted.

If $e=(c_0,c_1,c_2,c_3,-c_4)$, changing $Y$ to $-Y$ gives
\[
 J(c,e)=|\kappa_0|^2-\epsilon\sum_{j=1}^3|\kappa_j|^2
                    +\epsilon|\kappa_4|^2.
\]
Both shared quadratic penalties remain. The numerical engine first optimizes
this paired family and then performs unrestricted alternating best responses
within the five-mode image class. These searches do not certify their suprema.

\section*{2. An exact one-point scalar state}
There are independent variables $Z,G\sim\CN(0,1)$,
$R\sim\Gamma(m,1)$ and $V\sim\Gamma(d-1,1)$ such that the entire tuple
\eqref{eq:five-modes} has the law obtained from
\begin{equation}\label{eq:scalar-state}
 q=\frac{RG}{\sqrt D},\qquad
 A_X=\frac{R-m}{\sqrt m},\qquad
 A_Y=\frac{|G|^2+V-d}{\sqrt d}.
\end{equation}
To prove this, condition on $X$ and project the proper Gaussian $Y$ onto the unit
symmetric-square direction of $X$. Its projection is a standard proper Gaussian
independent of the orthogonal residual, whose squared norm is gamma of shape
$d-1$. The conditional distribution does not depend on $X$'s direction; the
squared projection coefficient is $R^2/D$. This proves the joint statement,
including $\|Y\|^2=|G|^2+V$.

The scalar $G$ in \eqref{eq:scalar-state} is not a new first-chaos coordinate.
It is a distributional coordinate of the original cubic. In particular $A_Y$
must not be replaced by an independent neutral Gaussian. The exact correlations
are
\[
 \E[A_X|q|^2]=\frac2{\sqrt m},\qquad
 \E[A_Y|q|^2]=\frac1{\sqrt d},\qquad \E A_XA_Y=0.
\]
More generally, for integer $k\ge1$,
\[
 \E[A_X|q|^{2k}]=\frac{2k}{\sqrt m}\E|q|^{2k},\qquad
 \E[A_Y|q|^{2k}]=\frac{k}{\sqrt d}\E|q|^{2k}.
\]
These follow from gamma integration by parts; their moment versions, through
$k=4$, and the complete five-mode Gram are checked over exact rationals in the
standard-library verifier. An independent $A_Y$ surrogate fails even the first
mixed moment.

For a smooth radial test $F$ for which the integrals and integration by parts
are justified, the stronger identities are
\[
 \E[A_XF(|q|^2)]=\frac2{\sqrt m}\E[|q|^2F'(|q|^2)],\qquad
 \E[A_YF(|q|^2)]=\frac1{\sqrt d}\E[|q|^2F'(|q|^2)].
\]
They apply, for example, to compactly supported $C^1$ tests. In an ansatz whose
visible conditional mean depends only on $|q|^2$, the two shared quadratic
coefficients are therefore tied. They cannot be chosen independently to
saturate separate projection budgets. This extra relation is not asserted for
general responses depending separately on $A_X,A_Y$.

There is an order-one spectral error if that warning is ignored. Keep the true
radius pair $(R,R')$ under original $X$ noise, but independently give the scalar
$G$ first-chaos correlation $\rho$. Since
$\E[RR']=m^2+m\rho^2$, this artificial model gives
\[
 \E[q\bar q']_{\rm artificial}=\frac{m\rho+\rho^3}{m+1},\qquad
 \E[q\bar q']_{\rm original}=\rho^3.
\]
Their difference is exactly $m\rho(1-\rho^2)/(m+1)>0$ for
$m>0$ and $0<\rho<1$. Thus convergence of a one-point law to a Gaussian does
not turn a third-chaos limit into a first-chaos noise coordinate. The factorized
gap and its strictness are Lean theorems. The covariance formulas follow from
the original Gaussian moments, not from those algebra theorems alone.

\section*{3. Two replicas with the original Gaussian coupling}
A one-point state does not determine a Hermite spectrum. For real
$0\le\rho\le1$, couple the original coordinates by
$X'=\rho X+\sqrt{1-\rho^2}\widetilde X$, and likewise for $Y,Z$, with fresh
independent original Gaussian blocks. Let
\[
 R=\|X\|^2,\quad R'=\|X'\|^2,\quad C=\sum_iX_i\bar X'_i.
\]
Their Gram is a two-by-two complex Wishart state. The symmetric-square probe
Gram has entries $R^2,R'^2,C^2$. In a common unitary basis of the two probe
directions, choose coefficients
\[
 a_1=\frac{\bar C^2}{R},\qquad
 a_2=\sqrt{R'^2-|C|^4/R^2}.
\]
Cauchy--Schwarz makes the radicand nonnegative. If
$(y_1,y_2)$ and $(y'_1,y'_2)$ are proper Gaussian coordinate pairs with
$\E[y_j\bar y'_k]=\rho\delta_{jk}$, then
\[
 q=R\bar y_1/\sqrt D,\qquad
 q'=(a_1\bar y'_1+a_2\bar y'_2)/\sqrt D.
\]
In particular
\[
 \E[q\bar q'\mid X,X']=\rho C^2/D,\qquad
 \E[q\bar q']=\rho^3.
\]
The common unitary basis depends on $X,X'$, but the conditional Gaussian
distribution of the jointly isotropic $Y,Y'$ does not. The remaining squared
norms form an independent two-by-two Wishart state of dimension $d-2$.
Consequently
\[
 \|Y\|^2=|y_1|^2+|y_2|^2+V_0,\qquad
 \|Y'\|^2=|y'_1|^2+|y'_2|^2+V_1,
\]
with $(V_0,V_1)$ the diagonal of that residual Wishart state. This retains the
dependence of both quadratic drivers on the same cubic projection coordinates.

Each Wishart state can be sampled from two gamma variables and one proper
complex Gaussian: for dimension $k$, take $r\sim\Gamma(k,1)$,
$s\sim\Gamma(k-1,1)$ and $h\sim\CN(0,1)$ independently, and put
\[
 C=\rho r+\sqrt{1-\rho^2}\sqrt r\bar h,
\quad
 r'=\rho^2r+(1-\rho^2)(|h|^2+s)
       +2\rho\sqrt{1-\rho^2}\sqrt r\Re h.
\]
At $k=1$, set $s=0$. Together with the visible pair and two projected $Y$
coordinate pairs, this uses twenty real scalar uniforms, independent of $m$.
It is a reparameterization of the original two-replica law, not an independent
Gaussian replacement. Endpoint $\rho=1$ is interpreted with zero new noise.

For any coefficient pair $c,e$ this state determines the actual kernel
\[
 K_{c,e}(\rho)=\E[\overline{f_c(Z,X,Y)}f_e(Z',X',Y')].
\]
Thus lower-witness adversaries can be carried into upper-pool discovery with
their original correlation structure. A certified full upper score still
requires coefficient enclosures and an infinite-tail bound. Neither follows
from these sampling formulas or the five-mode low-degree calculation alone.

\section*{4. A full-tail screen without fitting tail coefficients}
For charge-one disc-valued responses let
$c_k=\Re\langle P_{k+1,k}f,P_{k+1,k}g\rangle$. Parseval and
Cauchy--Schwarz give $\sum_k|c_k|\le1$. Taking the real part of the original
Gaussian noise kernel gives
\[
 K(\rho)=\sum_{k\ge0}c_k\rho^{2k+1}.
\]
Real coefficients can also be realized by averaging the response pair with
its conjugate-reflected pair, so this is a legitimate response-mixture score.
Put $T=\sum_{k\ge2}|c_k|$ and
$R_i=K(\rho_i)-c_0\rho_i-c_1\rho_i^3$, with $0<\rho_i\le1$.
Then $|R_i|\le\rho_i^5 T$, so the complete rounding score obeys
\[
 D(c)=c_0-|c_1|-T
 \le c_0-|c_1|-\frac{|R_i|}{\rho_i^5}.
\]
This is an upper screen on a candidate's achievable score, not a certified
lower score that would prove a Grothendieck upper bound.

More usefully, choose $|s|\le1$ and real prices $a_i$ satisfying
$\sum_i|a_i|\rho_i^5\le1$. Then
\begin{equation}\label{eq:pool-price}
 D(c)\le c_0-sc_1-\sum_i a_iR_i.
\end{equation}
Indeed $sc_1\le|c_1|$ and
$\sum_i a_iR_i\le\sum_i|a_i|\rho_i^5T\le T$.
The right side is affine in a response column. It therefore bounds every
convex mixture of a fixed pool by the largest priced column. Minimizing that
largest value is a small linear program on physical kernel samples. An
optimized upper screen below the incumbent score would exclude that pool,
provided the kernel integrals and low coefficients were rigorously enclosed.
It does not exclude a larger pool or give a universal witness cap.

The same prices define an admissible lower-witness multiplier sequence:
\[
 \lambda_0=1,\qquad\lambda_1=-s,\qquad
 \lambda_k=-\sum_i a_i\rho_i^{2k+1}\quad(k\ge2).
\]
The budget gives $|\lambda_k|\le1$ at every index; this all-index statement is
also proved in Lean. Its payoff on a response column is exactly
\eqref{eq:pool-price}. Thus a failed upper pool produces a concrete witness
against which to search for a missing response. A low payoff on that finite
pool does not bound the witness's global norm. Any new adversary must return
to the pool with its original-noise kernel, rather than merely its cubic score.
Equivalently, on charge-one functions its operator is
\[
 T_a=\left(1+\sum_i a_i\rho_i\right)P_{1,0}
 +\left(-s+\sum_i a_i\rho_i^3\right)P_{2,1}
 -\sum_i a_i\mathsf T_{\rho_i},
\]
where $\mathsf T_\rho$ is the original Gaussian noise operator and
$\mathsf T_1=I$. Thus the complete multiplier sequence is represented by a
finite collection of physical noise channels. No global cap or finite-matrix
transfer certificate for $T_a$ is asserted here.
Unlike the cubic operator in Section~1, $T_a$ need not have a five-mode image:
noise applied to a phase produces higher-degree components. Searching the
five-mode phase family against it is a restricted adversary search, not an
exact best-response oracle for the new operator.

There is a direct error ledger. If $c_0,c_1,K(\rho_i)$ are enclosed about
midpoints with errors $e_0,e_1,e_i$, add
\[
 e_0+|s|e_1+\sum_i|a_i|(e_i+\rho_i e_0+\rho_i^3e_1)
\]
to the midpoint right side of \eqref{eq:pool-price} for each column. This
supports interval certification without recovering millions of individual
chaos coefficients. Sampling variation is not such an enclosure.

\section*{5. An exact obstruction to the first price-induced witness}
The numerical adversary search suggests a simpler analytic baseline.
Let $f_m=\ph(Z+q_m/\sqrt3)$ and $g_m=\ph(Z-q_m/\sqrt3)$.
For every fixed real noise correlation, the replica construction gives
\[
 (q_m,q'_m)\ \Longrightarrow\ (G,G'),\qquad
 \E[G\bar G']=\rho^3,
\]
where the limit is a proper Gaussian pair independent of the visible pair.
Indeed $R/m,R'/m,C/m$ tend in probability to $1,1,\rho$ by the Gaussian law
of large numbers. The normalized projection coefficients tend to
$1,\rho^2,\sqrt{1-\rho^4}$ in the common two-coordinate representation.
This proves the joint limit directly, with its original noise grade retained.
The phase arguments have continuous nondegenerate Gaussian limits, so bounded
convergence under a coupling, or the bounded continuous mapping theorem,
gives convergence of their kernels.

Write the Gaussian phase kernel as
\[
 H(t)=\frac\pi4\sum_{k\ge0}
   \frac{(1/2)_k^2}{k!(k+1)!}t^{2k+1}.
\]
The limiting response-pair kernel is exactly
\[
 K_\infty(\rho)=H\left(\frac{3\rho-\rho^3}{4}\right),\qquad
 c_0=\frac{3\pi}{16},\quad c_1=-\frac{101\pi}{2048}.
\]
This is a classical cubic preprocessor of the Gaussian phase kernel. The
new search has therefore rediscovered a known type of rounding ingredient;
it should be included as an analytic baseline before further pool optimization.

The low coefficients also converge. This requires more than pointwise kernel
convergence alone. Here the uniform Parseval bound $\sum_k|c_k^{(m)}|\le1$
gives
\[
 |K_m(r)/r-c_0^{(m)}|\le r^2,\qquad
 \left|\frac{K_m(r)-c_0^{(m)}r}{r^3}-c_1^{(m)}\right|\le r^2.
\]
First fix $r$, pass to the kernel limit, and then let $r$ decrease to zero;
the first estimate gives convergence of $c_0^{(m)}$, and the second gives
convergence of $c_1^{(m)}$. Thus every finite priced payoff converges as well.

There is a qualification when treating such a limit as an upper-pool column:
componentwise coefficient convergence alone need not give convergence of the
absolute tail. Apply $\mathsf T_{\sqrt\eta}$ to both finite responses first,
with $0<\eta<1$. The mixed coefficients become
$\eta^{2k+1}c_k^{(m)}$. The same finite-head extraction works at every fixed
index; together with the uniform absolute-mass bound it gives $\ell^1$
convergence for fixed $\eta$, by splitting a finite head and a geometric tail.
Then let $\eta\uparrow1$ for the summable limiting sequence. A diagonal choice
proves membership in the $\ell^1$ closure of actual disc-valued response columns.
This does not assert a uniform tail for unsmoothed cubic phases or provide a
quantitative finite dimension here.

For the rationalized nine-column price saved with the replay data, an exact
standard-library series certificate gives limiting payoff
\[
 0.7169925532763968\ldots>\frac{7169}{10000}.
\]
It uses Machin's identity with alternating rational arctangent remainders to
enclose $\pi$. At each required argument $0\le t\le1/2$, the displayed
Haagerup series has successive term ratio at most $t^2\le1/4$, so a geometric
remainder encloses its entire tail. Forty-eight terms give an enclosure of
width less than $10^{-25}$. The interval verifier recomputes all terms and
the exact price budget; no quadrature or stochastic estimate is used.
Consequently the global norm of \emph{that price-induced witness} exceeds
$0.7169$, already on limits of genuine finite cubic responses, and sufficiently
large finite $m$ have the same strict obstruction. This rejects that witness
as an improvement on the incumbent cap. It concerns neither the incumbent
cubic witness nor the original-response conjecture.

\paragraph{Formal and numerical controls.}
The projection coefficient, covariance identity, nonnegative residual and
variance reconstruction are proved in
\texttt{Grothe/LowerAnalytic/ReplicaState.lean}, together with the artificial
noise gap and strictness. The tail sum, single-sample full-score bound and
finite priced screen are proved in
\texttt{Grothe/Upper/KernelScreen.lean}, as is all-index admissibility of the
price-induced witness. The invariant classification
and Wishart conditioning are the paper arguments above. Exact one-point mixed
moments and floating replica cross-Gram controls are separate checks. The exact
cross-Gram target is
$\operatorname{diag}(\rho,\rho^3,\rho^3,\rho^3,\rho^3)$.
Floating checks and optimizations do not certify a norm or an improved endpoint.

\end{document}
